% page 437 du fac-simile (image p-436 du scan)
% bloc 162.6 x 246.3 mm ; marge gauche 39.4 mm
\pgc{17.780mm}{0.000mm}{
\l{\cel{49}429.}
\l{}
\l{\sou{pedikul}o.I. (\sou{pato}l.) Parto streta qua suportas ula tumori.}
\l{\cel{3}II. (\sou{historio naturala}). Suportilo longa, gracila (che la}
\l{\cel{3}fungo, likeno, ovulo, e c.) EFIS.}
\l{}
\l{pedikuloso.\cel{2}(patol.)\cel{2}Ensemblo de la manifesti morbala, pro-}
\l{\cel{3}duktita da lausi sur korpo-parto, o mem sur la tota korpo,}
\l{\cel{3}per granda quanto de lausi. - DEFIS.}
\l{}
\l{\sou{pedonklo}. I. (\sou{bot.)}\cel{3}Suportilo di floro. - Kaudo di frukto}
\l{\cel{3}(kom ex.: che cerizo). - II. (\sou{zool}.)Prolonguro di la bulbo}
\l{\cel{3}spinala. - EFIS.}
\l{}
\l{\sou{pego.}\cel{3}(\sou{zool.}) Ucelo ek la klaso \textquotedbl{}klimeri\textquotedbl{}, qua havas beko}
\l{\cel{3}rekta, anguloza, e qua, per perforar la kortico di la arbori,}
\l{\cel{3}kaptas e devoras la insekti qui minas ta arbori. -\cel{2}L.\cel{1}\sou{picus}.}
\l{}
\l{\sou{pegazo}. I. (mitol.) Kavalo aloza q ua spricigis de la monto}
\l{\cel{3}Helikon (Grekia) la fonto di Hippokrenes ube onu cherpis}
\l{\cel{3}la inspirado poeziala. - II. (\sou{metaf.}) Simbolo di la inspirado}
\l{\cel{3}poeziala. - DEFIRS.}
\l{}
\l{\sou{peizajo}. I. Parto di lando, plu o min vasta, e plu o min pikt-}
\l{\cel{3}inda. - II. (\sou{metaf}.)Pikturi qua reprezentas tala parto.}
\l{\cel{3}FIRS.}
\l{}
\l{\sou{pejorativa}. (\sou{gram.)}\cel{2}Qua donas senco desfavoroza, desprizigiva.}
\l{\cel{3}- DEFI.}
\l{}
\l{\sou{pekar}. (\sou{netrans.})\cel{2}(\sou{religio krist.}) Kulpar kontre la lego}
\l{\cel{3}de Deo, kontre la lego religiala. - EFIS.}
\l{}
\l{\sou{pekario}. (\sou{zool.}) Mamifero di Amerika, kelke simila a porko. -}
\l{\cel{3}L.\cel{1}\sou{dicotylus torquatus}.\cel{2}DEFIS.}
\l{}
\l{\sou{pekino}. Stofo de silko en qua strio satinatra alternas kun strio}
\l{\cel{3}matida. -}
\l{}
\l{\sou{pektar}. (\sou{trans.)}\cel{4}Desintrikar, netigar, akomodar (sua hararo,}
\l{\cel{3}sua barbo) per instrumento qua havas denti de materio korna,}
\l{\cel{3}skalia, ivora, metala. - FIS.}
\l{}
\l{\sou{pektino}. (\sou{biol.})\cel{2}Substanco neutra, quan kontenas la frukti}
\l{\cel{3}matura. - DEFIRS.}
\l{}
\l{\sou{pektoro}. (\sou{anat.)}\cel{2}Parto di la korpo,qua kontenas la pulmoni e}
\l{\cel{3}la kordio. - EFIS.}
\l{}
\l{\sou{pektoso.}\cel{1}(\sou{biol.})\cel{2}Kompozajo qua naskas en la tisui vejetantala,}
\l{\cel{3}qua kontributas (kun kutoso e vaskuloso) ad aglutinar la}
\l{\cel{3}celuloso-fibri. - DEFIRS.}
\l{}
\l{\sou{pekunio.}\cel{1}Omna speco di moneto (metala, papera, e c.) - DEFIRS.}
\l{}
\l{pelo. (\sou{ana}t.)\cel{2}Membrano extera qua tegas omna parti di la korpo}
\l{\cel{3}di homo o di animalo. - DFIS.}
}
